\documentclass[10pt,a4paper]{amsart}
\usepackage{amsmath,amssymb,amsthm,graphicx}
\usepackage{verbatim}
\newenvironment{reference}{\comment}{\endcomment}
%\newenvironment{reference}{}{}
\newenvironment{slogan}{\comment}{\endcomment}
\newenvironment{history}{\comment}{\endcomment}

% For commutative diagrams we use Xy-pic
\usepackage[all]{xy}

% We use 2cell for 2-commutative diagrams.
\xyoption{2cell}
% Theorem environments.
%
\theoremstyle{plain}
\newtheorem{theorem}[subsection]{Theorem}
\newtheorem{proposition}[subsection]{Proposition}
\newtheorem{lemma}[subsection]{Lemma}

\theoremstyle{definition}
\newtheorem{definition}[subsection]{Definition}
\newtheorem{example}[subsection]{Example}
\newtheorem{exercise}[subsection]{Exercise}
\newtheorem{situation}[subsection]{Situation}

\theoremstyle{remark}
\newtheorem{remark}[subsection]{Remark}
\newtheorem{remarks}[subsection]{Remarks}

\numberwithin{equation}{subsection}

% Macros
%
\def\lim{\mathop{\mathrm{lim}}\nolimits}
\def\colim{\mathop{\mathrm{colim}}\nolimits}
\def\Spec{\mathop{\mathrm{Spec}}}
\def\Hom{\mathop{\mathrm{Hom}}\nolimits}
\def\Ext{\mathop{\mathrm{Ext}}\nolimits}
\def\SheafHom{\mathop{\mathcal{H}\!\mathit{om}}\nolimits}
\def\SheafExt{\mathop{\mathcal{E}\!\mathit{xt}}\nolimits}
\def\Sch{\mathit{Sch}}
\def\Mor{\mathop{\mathrm{Mor}}\nolimits}
\def\Ob{\mathop{\mathrm{Ob}}\nolimits}
\def\Sh{\mathop{\mathit{Sh}}\nolimits}
\def\NL{\mathop{N\!L}\nolimits}
\def\CH{\mathop{\mathrm{CH}}\nolimits}
\def\proetale{{pro\text{-}\acute{e}tale}}
\def\etale{{\acute{e}tale}}
\def\QCoh{\mathit{QCoh}}
\def\Ker{\mathop{\mathrm{Ker}}}
\def\Im{\mathop{\mathrm{Im}}}
\def\Coker{\mathop{\mathrm{Coker}}}
\def\Coim{\mathop{\mathrm{Coim}}}

% Boxtimes
%
\DeclareMathSymbol{\boxtimes}{\mathbin}{AMSa}{"02}

%
% Macros for moduli stacks/spaces
%
\def\QCohstack{\mathcal{QC}\!\mathit{oh}}
\def\Cohstack{\mathcal{C}\!\mathit{oh}}
\def\Spacesstack{\mathcal{S}\!\mathit{paces}}
\def\Quotfunctor{\mathrm{Quot}}
\def\Hilbfunctor{\mathrm{Hilb}}
\def\Curvesstack{\mathcal{C}\!\mathit{urves}}
\def\Polarizedstack{\mathcal{P}\!\mathit{olarized}}
\def\Complexesstack{\mathcal{C}\!\mathit{omplexes}}
% \Pic is the operator that assigns to X its picard group, usage \Pic(X)
% \Picardstack_{X/B} denotes the Picard stack of X over B
% \Picardfunctor_{X/B} denotes the Picard functor of X over B
\def\Pic{\mathop{\mathrm{Pic}}\nolimits}
\def\Picardstack{\mathcal{P}\!\mathit{ic}}
\def\Picardfunctor{\mathrm{Pic}}
\def\Deformationcategory{\mathcal{D}\!\mathit{ef}}
\title{Existence and uniqueness of the total Chern class in a graded cohomology theory}
\author{The Stacks Project Authors}
\date{stacks-project:weil.tex:local-snapshot-20261012; GNU FDL 1.2 or later}
\begin{document}
\maketitle
\noindent\textit{Editorial note.} All delivered mathematical text is the author's own, copied from the pinned Stacks Project chapter file weil.tex; no statement or proof step was rewritten, reordered or removed.\par\medskip
\noindent\textit{Editorial note.} The counted proof is the author's complete argument: the construction of the Chern classes through the projective bundle relation, the compatibility with pullbacks, the multiplicativity in short exact sequences by induction on the rank, and the uniqueness argument.\par\medskip
\noindent\textit{Editorial note.} The label proposition-chern-class names the counted statement, and the proof environment immediately following it in the source is delivered verbatim as its proof.\par\medskip
\noindent\textit{Editorial note.} The counted statement is carried by the author's own proposition environment in the chapter source.\par\medskip
\noindent\textit{Editorial note.} Pre-open gate: python operations/agent-authoring/tools/precard\_gate.py --source weil.tex --label proposition-chern-class --start 4124 --end 4283 --local-id 28136 returned RESULT=OK (L28136 free, occupancy FREE, public interval clear). Pre-assembly tag self-check: canonical tag resolved and not present in reference/public-items.jsonl.\par\medskip
\noindent\textit{Editorial note.} The statement and proof are copied from the pinned chapter file; the cited inputs are kept as external pointers with their canonical Stacks tags.\par\medskip
\noindent\textit{Allowed standard prerequisites (editorial).}\begin{itemize}
\item The projective bundle formula for the graded cohomology theory. The standard description of the cohomology of the projective bundle of a finite locally free module, giving the relation used to define the Chern classes. Supplies the relation in the cohomology of the projective bundle from which the Chern classes of a finite locally free module are extracted.
\item Cohomological vanishing for a zero scheme of a regular section. The axiom of the cohomology theory describing the effect of the class of a zero scheme of a regular section cut out by an invertible module. Turns the vanishing of the pulled back defining relation on the projective bundle of the quotient into the identity expressing the multiplicativity of the total Chern class.
\item Ampness, very ampleness and Bertini type statements for the construction of the Chern classes. The construction of the projective bundle of a finite locally free module and its functoriality for the compatibility with pullbacks. Provides the functoriality of the projective bundle used to compare the defining relations before and after a pullback.
\end{itemize}\medskip
\begin{proposition}
\label{proposition-chern-class}
In the situation above there is a unique rule which assigns to
every $X \in \Ob(\mathcal{C})$ a ``total Chern class''
$$
c^A : K_0(\textit{Vect}(X)) \longrightarrow  \prod\nolimits_{i \geq 0} A^i(X)
$$
with the following properties
\begin{enumerate}
\item For $X \in \Ob(\mathcal{C})$ we have
$c^A(\alpha + \beta) = c^A(\alpha) c^A(\beta)$
and $c^A(0) = 1$.
\item If $f : X' \to X$ is a morphism of $\mathcal{C}$, then
$f^* \circ c^A =  c^A \circ f^*$.
\item Given $X \in \Ob(\mathcal{C})$ and $\mathcal{L} \in \Pic(X)$
we have $c^A([\mathcal{L}]) = 1 + c_1^A(\mathcal{L})$.
\end{enumerate}
\end{proposition}

\begin{proof}
Let $X \in \Ob(\mathcal{C})$ and let $\mathcal{E}$ be a finite
locally free $\mathcal{O}_X$-module. We first show how to define
an element $c^A(\mathcal{E}) \in A(X)$.

\medskip\noindent
As a first step, let $X = \bigcup X_r$ be the decomposition into
open and closed subschemes such that $\mathcal{E}|_{X_r}$ has
constant rank $r$. Since $X$ is quasi-compact, this decomposition
is finite. Hence $A(X) = \prod A(X_r)$. Thus it suffices to define
$c^A(\mathcal{E})$ when $\mathcal{E}$ has constant rank $r$. In this
case let $p : P \to X$ be the projective bundle of $\mathcal{E}$.
We can uniquely define elements $c_i^A(\mathcal{E}) \in A^i(X)$
for $i \geq 0$ such that $c_0^A(\mathcal{E}) = 1$ and the equation
\begin{equation}
\label{equation-chern-classes}
\sum\nolimits_{i = 0}^r
(-1)^i c_1(\mathcal{O}_P(1))^i \cup p^*c^A_{r - i}(\mathcal{E})
= 0
\end{equation}
is true. As usual we set
$c^A(\mathcal{E}) = c_0^A(\mathcal{E}) + c_1^A(\mathcal{E}) + \ldots
+ c_r^A(\mathcal{E})$ in $A(X)$.

\medskip\noindent
If $\mathcal{E}$ is invertible, then
$c^A(\mathcal{E}) = 1 + c_1^A(\mathcal{L})$.
This follows immediately from the construction above.

\medskip\noindent
The elements $c_i^A(\mathcal{E})$ are in the center of $A(X)$.
Namely, to prove this we may assume $\mathcal{E}$ has constant rank $r$.
Let $p : P \to X$ be the corresponding projective bundle.
if $a \in A(X)$ then $p^*a \cup (-1)^r c_1(\mathcal{O}_P(1))^r =
(-1)^r c_1(\mathcal{O}_P(1))^r \cup p^*a$ and hence we must have the same
for all the other terms in the expression defining $c_i^A(\mathcal{E})$
as well and we conclude.

\medskip\noindent
If $f : X' \to X$ is a morphism of $\mathcal{C}$, then
$f^*c_i^A(\mathcal{E}) = c_i^A(f^*\mathcal{E})$.
Namely, to prove this we may assume $\mathcal{E}$ has constant rank $r$.
Let $p : P \to X$ and $p' : P' \to X'$ be the projective
bundles corresponding to $\mathcal{E}$ and $f^*\mathcal{E}$.
The induced morphism $g : P' \to P$ is a morphism of $\mathcal{C}$.
The pullback by $g$ of the equality defining $c_i^A(\mathcal{E})$
is the corresponding equation for $f^*\mathcal{E}$ and we conclude.

\medskip\noindent
Let $X \in \Ob(\mathcal{C})$. Consider a short exact sequence
$$
0 \to \mathcal{L} \to \mathcal{E} \to \mathcal{F} \to 0
$$
of finite locally free $\mathcal{O}_X$-modules with $\mathcal{L}$ invertible.
Then
$$
c^A(\mathcal{E}) = c^A(\mathcal{L}) c^A(\mathcal{F})
$$
Namely, by the construction of $c^A_i$ we may assume $\mathcal{E}$ has
constant rank $r + 1$ and $\mathcal{F}$ has constant rank $r$.
The inclusion
$$
i : P' = \mathbf{P}(\mathcal{F}) \longrightarrow \mathbf{P}(\mathcal{E}) = P
$$
is a morphism of $\mathcal{C}$ and it is the zero scheme of a regular
section of the invertible module
$\mathcal{L}^{\otimes -1} \otimes \mathcal{O}_P(1)$.
The element
$$
\sum\nolimits_{i = 0}^r (-1)^i c_1^A(\mathcal{O}_P(1))^i \cup
p^*c^A_i(\mathcal{F})
$$
pulls back to zero on $P'$ by definition. Hence we see that
$$
\left(c_1^A(\mathcal{O}_P(1)) - c_1^A(\mathcal{L})\right) \cup
\left(\sum\nolimits_{i = 0}^r (-1)^i c_1^A(\mathcal{O}_P(1))^i \cup
p^*c^A_i(\mathcal{F})\right) = 0
$$
in $A^*(P)$ by assumption (5) on our cohomology $A$.
By definition of $c_1^A(\mathcal{E})$
this gives the desired equality.

\medskip\noindent
Let $X \in \Ob(\mathcal{C})$. Consider a short exact sequence
$$
0 \to \mathcal{E} \to \mathcal{F} \to \mathcal{G} \to 0
$$
of finite locally free $\mathcal{O}_X$-modules. Then
$$
c^A(\mathcal{F}) = c^A(\mathcal{E}) c^A(\mathcal{G})
$$
Namely, by the construction of $c^A_i$ we may assume
$\mathcal{E}$, $\mathcal{F}$, and $\mathcal{G}$ have
constant ranks $r$, $s$, and $t$. We prove it by induction on $r$.
The case $r = 1$ was done above. If $r > 1$, then it suffices to check
this after pulling back by the morphism $\mathbf{P}(\mathcal{E}^\vee) \to X$.
Thus we may assume we have an invertible submodule
$\mathcal{L} \subset \mathcal{E}$ such that both
$\mathcal{E}' = \mathcal{E}/\mathcal{L}$ and
$\mathcal{F}' = \mathcal{E}/\mathcal{L}$ are finite locally free
(of ranks $s - 1$ and $t - 1$). Then we have
$$
c^A(\mathcal{E}) = c^A(\mathcal{L}) c^A(\mathcal{E}')
\quad\text{and}\quad
c^A(\mathcal{F}) = c^A(\mathcal{L}) c^A(\mathcal{F}')
$$
Since we have the short exact sequence
$$
0 \to \mathcal{E}' \to \mathcal{F}' \to \mathcal{G} \to 0
$$
we see by induction hypothesis that
$$
c^A(\mathcal{F}') = c^A(\mathcal{E}') c^A(\mathcal{G})
$$
Thus the result follows from a formal calculation.

\medskip\noindent
At this point for $X \in \Ob(\mathcal{C})$
we can define $c^A : K_0(\textit{Vect}(X)) \to A(X)$.
Namely, we send a generator $[\mathcal{E}]$ to $c^A(\mathcal{E})$
and we extend multiplicatively. Thus for example
$c^A(-[\mathcal{E}]) = c^A(\mathcal{E})^{-1}$ is the formal
inverse of $a^A([\mathcal{E}])$.
The multiplicativity in short exact sequences shown above
guarantees that this works.

\medskip\noindent
Uniqueness. Suppose $X \in \Ob(\mathcal{C})$ and $\mathcal{E}$
is a finite locally free $\mathcal{O}_X$-module. We want to show
that conditions (1), (2), and (3) of the lemma uniquely determine
$c^A([\mathcal{E}])$. To prove this we may assume $\mathcal{E}$
has constant rank $r$; this already uses (2). Then we may use induction on $r$.
If $r = 1$, then uniqueness follows from (3).
If $r > 1$ we pullback using (2) to the projective bundle $p : P \to X$
and we see that we may assume we have a short exact sequence
$0 \to \mathcal{E}' \to \mathcal{E} \to \mathcal{E}'' \to 0$
with $\mathcal{E}'$ and $\mathcal{E}''$ having lower rank.
By induction hypothesis $c^A(\mathcal{E}')$ and $c^A(\mathcal{E}'')$
are uniquely determined. Thus uniqueness for $\mathcal{E}$ by
the axiom (1).
\end{proof}

\end{document}
